\section{Maurras and Dumezil}

Early in his career Dumezil was involved with Charles Maurras and the Action francaise. In his biography \textit{Maurras: Le chaos et l'ordre}, Stephane Giocanti describes their relationship and the influence that Maurras had over Dumezil. A translation of the relevant passage follows:

\begin{quotex}
... in the 20s, and up until its final decline, Action Française regularly attracted new spirits, notably among students: Jacques Isorni, Tixier-Vignancour. Around 1925, Pierre Gaxotte introduced Maurras to George Dumezil, who would have a huge career as an historian and anthropologist. Maurras took them both on as secretaries, one for the day shift, the other for the night.

Author of a doctoral thesis on ``The Banquet of Immortality'', Dumezil wrote for \emph{La Revue universelle} and \emph{L'Action française}, since 1920, notably articles on ``monarchical and dynastic principles'' which were not without agreement with his later thought. His brief correspondence with Maurras attests to a deferential and friendly relationship. In May 1925, he told Maurras about his engagement. ``I will go to see you in two days, after a last stop in Dijon. But I am eager to inform you right away because on Saturday or Sunday I am leaving for Prague and I would not forgive myself if I caused any neglect in the office. Maurras sent his wishes for a good trip. In September, while Dumezil was in Turkey where he wanted the chair of the history of religions, Maurras had sent him his latest books: ``My dear master, your books, your two books were waiting for me at the Sarmatians. ... They are going to shake up the sweet barbaric dust and give me the taste and the strength for battle. You satisfy me with healthy favours that I will try some day to recognize. ... ''

Beyond the anecdotal aspects of this relation, we can ask about Dumezil's debt to Maurras. We know that the historian would elaborate some comparative methods to throw light on the depths of the human spirit, by exposing their laws, structure, and hierarchies. He had a scientific spirit. Maurras, in developing the birth and the triumph of reason in classical Greece and in showing its centrality in European civilisation, was a poet-archaeologist of the human spirit. Instead of proceeding scientifically, \emph{Anthinéa} proceeds by means of a blend of mythography, art, and philosophy. The reader only gets to a piece of philosophical anthropology through aesthetic, historic, and political meditations. Dumezil is therefore a son of \emph{Anthinéa}, \emph{a fortiori} when he asserts that the results of his research have an interest ``more aesthetic than scientific'' that is, far from materialism. In spite of the prudence this formidable erudite would observe after dealing with accusations, Maurras (not certainly the polemicist, but the thinker of the ancient world) can be considered as one of the intercessors of his thought and a founding authority in the formation of his anti-German and anti-Nazi prejudice.
\end{quotex}

\flrightit{Posted on 2008-04-10 by Cologero}

